\par\Needspace{40\baselineskip}\pdfbookmark[1]{Gauss to Schumacher. 1822 November 25.}{schumacher1822}\begin{center}\rule{0.72\linewidth}{0.3pt}\par Astronomische Nachrichten, edited by \textsc{Schumacher}. No. 24. Supplement 2. 1822 December.\par\rule{0.72\linewidth}{0.3pt}\end{center}

\begin{center}
\textsc{Gauss} to \textsc{Schumacher}.
\end{center}
\begin{flushright}
Göttingen, 1822 November 25.
\end{flushright}
--- --- I have often shared Professor \textsc{Encke}'s wish for logarithm tables with 6 digits, and you will certainly earn much gratitude if you arrange for such tables to be made. I endorse everything Mr \textsc{Encke} says about their outward form and contents as my own opinion, except that the proportional parts seem superfluous to me, and everything superfluous makes a table harder to use and take in at a glance. In such matters, much admittedly depends on individual habits; but if some users of tables have habits different from mine, others probably have the same habits. I therefore touch on one further point, namely the change in the \(4^{\mathrm{th}}\) digit in logarithm tables. You know the arrangement made in \textsc{Callet}'s tables in this respect, which some have regarded as an improvement. I admit that I hold the opposite opinion: for frequent use, I consider the \textit{regular} division into groups of 5 rows by horizontal lines every 5 rows, as in \textsc{Sherwin}'s and other tables, \textit{much} more essential and convenient, and I therefore never care to use \textsc{Callet}'s logarithms. In my many years of practice, I know of not a single instance in which using \textsc{Sherwin}'s tables led me into an error of calculation in the \(4^{\mathrm{th}}\) digit. I therefore attach no value at all to the asterisks in \textsc{Vega} and other tables, and prefer \textsc{Sherwin}'s because of their better paper and more attractive digits. Anyone who values such warning signs can easily add red or green dots at the relevant places in their own copy. ---
